
On
considère alors le morphisme canonique (2.7.2) :
@@M0@@
soit @@M1@@ le cône de ce morphisme. C'est encore un complexe @@M2@@-parfait : en effet, pour tout @@M3@@, @@M4@@ est un complexe parfait sur @@M5@@, car il en est ainsi quand @@M6@@ est borné à termes localement libres de type fini sur @@M7@@; par suite, le complexe @@M8@@ est @@M9@@-parfait, et il en est donc de même pour le cône (I 4.10).

Par ailleurs, on peut écrire sur tout ouvert @@M10@@ du recouvrement le diagramme commutatif :
@@M11@@
qui induit un isomorphisme entre les cônes; or l'homomorphisme du bas est un isomorphisme en degré @@M12@@; donc @@M13@@ est isomorphe sur @@M14@@ à un complexe borné à termes localement libres de type fini sur @@M15@@, dont les composantes @@M16@@ sont nulles pour @@M17@@. Par suite, @@M18@@ vérifie les mêmes propriétés que @@M19@@, @@M20@@ étant remplacé par @@M21@@. On peut donc itérer le processus, et il suffit de voir qu'au bout d'un nombre fini d'opérations, le complexe @@M22@@ obtenu est acyclique. Là encore la question est locale, si bien qu'il suffit de le voir pour @@M23@@. Or pour chaque terme @@M24@@, les morphismes considérés plus haut sont ceux qui définissent la filtration de la proposition 1.7; cette filtration est finie, donc en un nombre fini d'opérations, on arrive à partir de @@M25@@ au complexe nul, c'est-à-dire que le complexe @@M26@@ correspondant est acyclique.

